\section{The Rule Algebra}
\label{sec:rules}

\subsection{Definitions}

A rule $r$ has a compiled test $\Hit(r,q,t)$: whether $r$ selects the path $q$
whose final component has type $t$.  It is evaluated on folded text
(Definition~\ref{def:fold}) with a pattern that was folded the same way.

\begin{definition}[Exclusion]
\label{def:excl}
For a rule set $\mathcal R$ and a path $p$ of type $t$, let
$\Anc(p)=\{q\mid q\text{ is a prefix of }p\}$, where each proper prefix has type
$\Dir$ and $p$ has type $t$.  The last rule of $\mathcal R$ that hits $q$ decides
$q$: it excludes $q$ unless it is negated.  Then
\[
  \Match(\mathcal R,p,t)\;=\;\exists\, q\in\Anc(p):\ q\text{ is excluded by }\mathcal R.
\]
Because an excluded ancestor already makes the whole path excluded, a later
negation cannot re-include something inside it, as in git.
\end{definition}

\begin{definition}[Match below]
\label{def:matchbelow}
$\MatchBelow(\mathcal R,f,p,t)$ is $\Match$ restricted to the prefixes of $p$
that are strictly longer than $f$.  It is used to apply hide rules while listing
a directory: entering a hidden directory on purpose shows its contents.
\end{definition}

\subsection{Laws for a rule set without negation (Deny)}

\begin{law}[Deny monotonicity]
\label{law:deny-monotone}
If $\mathcal R\subseteq\mathcal R'$ then
$\Match(\mathcal R,p,t)\Rightarrow\Match(\mathcal R',p,t)$.  Adding a deny rule
never makes anything accessible.
\Tests{rules/laws_test.go::TestLawDenyIsMonotoneAndOrderIndependent}
\end{law}

\begin{law}[Order independence]
\label{law:deny-order}
If $\mathcal R'$ is a permutation of $\mathcal R$, then
$\Match(\mathcal R,p,t)=\Match(\mathcal R',p,t)$.
\Tests{rules/laws_test.go::TestLawDenyIsMonotoneAndOrderIndependent}
\end{law}

\begin{law}[No negation in deny]
\label{law:no-negation}
Parsing a deny file that contains a negated rule fails.  Together with
Laws~\ref{law:deny-monotone} and~\ref{law:deny-order} this gives the security
reading of a deny file: it is a plain \emph{set} of exclusions.
\Tests{rules/laws_test.go::TestLawDenyRefusesNegation}
\end{law}

\subsection{Laws for every rule set}

\begin{law}[Ancestor closure]
\label{law:ancestor}
If $\Match(\mathcal R,p,\Dir)$ then $\Match(\mathcal R,p\cdot n,t)$ for every
name $n$ and type $t$.  This holds also when $\mathcal R$ contains negated
rules.
\Tests{rules/laws_test.go::TestLawAncestorClosure}
\end{law}

\begin{law}[Fold invariance]
\label{law:fold-invariance}
If $\fold{p}=\fold{p'}$ then $\Match(\mathcal R,p,t)=\Match(\mathcal R,p',t)$:
a path is excluded exactly when any spelling that the file system treats as the
same name is.  In particular no spelling of a denied name (upper case, decomposed
accents, long s, Kelvin sign, sharp s, ligatures) escapes a rule.
\Tests{rules/laws_test.go::TestLawMatchDependsOnlyOnTheFoldedPath; guard/foldalias_test.go}
\end{law}

\begin{law}[Narrowing]
\label{law:narrowing}
$\MatchBelow(\mathcal R,f,p,t)\Rightarrow\Match(\mathcal R,p,t)$.  Ignoring
ancestors can only remove matches.
\Tests{rules/laws_test.go::TestLawMatchBelowNeverAddsMatches}
\end{law}

\begin{law}[Single component]
\label{law:component}
If a pattern contains neither \code{/} nor \code{**}, then $\Hit(r,q,t)$ depends
only on the final component of $q$.  Glob classes never match \code{/}: a
negated class excludes it, and a \code{/} inside a class is a pattern error
(fail closed).
\Tests{rules/laws_test.go::TestLawSlashlessPatternsStayWithinOneComponent; TestSlashInsideAClassIsAParseError}
\end{law}

\begin{law}[Credential folders match wherever they are]
\label{law:credential-anywhere}
The core rules for credential folders (\code{.ssh}, \code{.aws}, \code{.gnupg},
\code{.kube}, \code{.config/gh}, Keychains, and the browser profiles) and files
(\code{.netrc}) are not anchored to the current home directory: a path that
contains one of them at any depth is excluded.  The same secrets exist in a
backup or clone of the home folder, in another account's home, and under a wrong
\code{\$HOME}, where an anchored rule protects nothing.
\Tests{guard/differential\_test.go::TestCredentialFoldersAreDeniedWhereverTheyAre}
\end{law}

\begin{law}[A rule file works as written or fails loudly]
\label{law:rule-file}
For every rule file, either parsing fails with a message, or each rule line means
what a reader would take it to mean.  In particular no accepted line is a rule that
can never match: a leading byte order mark and any line ending are normalized, and
an inline comment, \code{\~{}user/}, a quoted pattern, an invisible character or a
lookalike of the slash are refused.
\Tests{rules/parse\_audit\_test.go}
\end{law}

\begin{law}[The quick test never changes an answer]
\label{law:fast-path}
Matching first asks whether any rule could match, using one combined expression
per kind of rule, and tests slash-free patterns against the last path component
only.  These are optimizations: for every rule set without negation and every
path, the result is that of consulting the rules one by one against the whole
path.
\Tests{rules/laws\_test.go::TestLawFastPathEqualsTheRulesOneByOne}
\end{law}

\begin{law}[Trailing whitespace is not significant]
\label{law:tail}
For every rule $r$ and path $q$, $\Hit(r,q\cdot\!\!\text{``}n\text{''},t)=
\Hit(r,q\cdot\!\!\text{``}n\ \text{''},t)$ for names whose trailing characters
are white space (space, tab, newline).  macOS permits such names and they must
not slip past a rule for \code{*.pem}.
\Tests{rules/rules_test.go; guard/review_test.go}
\end{law}

\subsection{Boundaries}

\begin{boundary}[Hard links]
Rules name paths, so a hard link to a denied file is invisible to them.  The guard
therefore judges a regular file with more than one name by its \emph{identity}
(device and inode): if any name of the file is denied, every name is refused, in
\code{Open} and in listings alike (Law~\ref{law:linked}).  The other names come
from an index of all multi-name files under the roots and under the protected
credential locations, built in the background because visiting every file takes
seconds.  A multi-name file is served only if the index accounts for all of its
names now: it found as many distinct names (by folder identity and folded name) as
the file has links, none is denied, and each is still, when the file is used, a real
path to that file.  Otherwise it is refused: while the index is not ready, for a
name outside every root and protected location or in a folder that could not be
read, after a walk stopped at its limit, and after the names changed.  Within one
listing or search a verdict is reused for at most one second, so a name or size
shown may lag a rename by that long; \code{Open} never reuses it.  Not detected:
a denied name outside every root and protected location, whose file is refused
rather than served.
\end{boundary}

\begin{boundary}[Case classes in ranges]
A positive class range such as \code{[+-0]} can include \code{/} by range even
though the class text does not.  No shipped rule uses such a range; a future
tightening would expand ranges before compiling.
\end{boundary}
